
\begin{proposition}[A rational ordinary-triple image has no positive-genus double branch cover]\label{cdp:st:C.2.D4.2}
Let $f:Y\to C_0$ have the torus-fibre and ambient hypotheses
of \ref{cdp:st:C.2.D4.1}. Let $H\subset S$ be a smooth rational curve
with generic ordinary-triple transversal. Suppose two of
the three generic normalization branches form a connected
degree-two cover $C\to H$, and the remaining branch forms
a degree-one cover. Then the smooth complete curve $C$
is rational. The assertion includes all zeros, with their
multiplicities, of the entire cubic normal tensor.

Consequently, in the irrational ruled setup of \ref{cdp:st:A.5.19},
a noncontracted coefficient-two horizontal conductor
with generic ordinary-triple transversal has smooth
elliptic image. Its normalization is smooth minimal
elliptic ruled, the horizontal cover is unramified of
degree three, and $e(S)=0$.
\end{proposition}
\begin{proof}
Suppose $g=g(C)\ge1$. Use the actual normal tensor and
normal-cone comparison (C.D4.1)--(C.D4.2), and write
\[
 N=N_{H/Y}=\OO(a)\oplus\OO(-2-a),\quad a\ge-1,
 \qquad P\in H^0(H,\operatorname{Sym}^3N^*),
 \qquad U=\{P=1\}\subset N.
\tag{C.DD.1}\label{cdp:eq:C.DD.1}
\]
The comparison embeds a compact core of $U$, including
its compact intersection pairing, into an actual nearby
torus. As proved there, $b_1(U)\le3$ and a nonzero
intersection pairing contradict this embedding.

Put $E=\PP(N)$, using lines in $N$, and let
$L=\OO_E(1)$, $\xi=c_1(L)$ and $F$ be the fibre class.
These conventions give
\[
 \xi^2=-\deg N=2,\qquad \xi F=1,\quad F^2=0,
 \qquad K_E=-2\xi.
\tag{C.DD.2}\label{cdp:eq:C.DD.2}
\]
Let $A$ be the common zero divisor of the coefficients
of $P$, and put $v=\deg A$. The divisor of the induced
section of $L^3$ is
\[
                    B_{\rm tot}=R_1+R_2+p^*A.
\tag{C.DD.3}\label{cdp:eq:C.DD.3}
\]
Here $R_1$ is a section and $R_2$ is an integral
degree-two curve over $H$. The section is smooth:
its rational tangent-root map extends over the smooth
complete base into the proper bundle $E$.
The normalization of $R_2$ is precisely $C$. Over the
generic point a normalization branch is identified
with its distinct tangent root; taking the integral
closure in this same quadratic function field identifies
the two complete normal covers, as in (C.RD.2).

\paragraph{All possible common divisors.}
For $a\ge0$ and $v>0$, the pure cube of the coordinate
in $\OO(a)$ has coefficient of degree $-3a-v$, and
therefore vanishes. The other coordinate is a factor.
The quadratic factor is irreducible over the function
field. Its coefficient of the first coordinate squared
must be nonzero; that coefficient has degree
\[
                              2-a-v.
\tag{C.DD.4}\label{cdp:eq:C.DD.4}
\]
Thus $a+v\le2$. In particular, $v\le2$, and $v=2$
forces $a=0$.

For $a=-1$, identify $E$ with the product of the base
and the projective direction line. Write $d\ge0$ for
the degree of the map defining the section $R_1$.
The bidegrees, with degree over $H$ first, are
\[
 R_1:(1,d),\qquad R_2:(2,3-v-d),\qquad
 p_a(R_2)=2-v-d.
\tag{C.DD.5}\label{cdp:eq:C.DD.5}
\]
The inequality $g\le p_a(R_2)$ gives $v+d\le1$.
Consequently the only common-divisor possibilities
for our positive-genus cover are $v=0$, $v=1$, and
$v=2$ with $N=\OO\oplus\OO(-2)$.

\paragraph{The exact cyclic-cover homology formulas.}
We record the comparison with relations and disconnected
boundary retained. First suppose that \eqref{cdp:eq:C.DD.3} is reduced,
and let $B=B_{\rm tot}$. The cyclic cover
\[
 \pi:X_0=\operatorname{Spec}_E
       (\OO_E\oplus L^{-1}\oplus L^{-2})\longrightarrow E
\tag{C.DD.6}\label{cdp:eq:C.DD.6}
\]
has multiplication given by $P$. It is normal, since
the branch divisor is reduced. Its complement over
$B$ is exactly $U$: if $w\in\ell^*$ is a nonzero
cube root of $P|_\ell$, its corresponding vector in
the line $\ell\subset N$ is the unique vector $z$
with $w(z)=1$. This description specifies inverse
maps and their bundle transitions.

All three summands in \eqref{cdp:eq:C.DD.6} have vanishing $H^1$.
Indeed the relative fibre cohomology of $L^{-1}$
vanishes, while relative duality gives
\[
 p_*L^{-2}=0,\qquad R^1p_*L^{-2}=\det N=\OO(-2).
\tag{C.DD.7}\label{cdp:eq:C.DD.7}
\]
Thus $H^1(E,L^{-2})=H^0(H,\OO(-2))=0$;
$H^1(E,\OO_E)=0$ as well. For a resolution
$\rho:X\to X_0$, including resolution of the reduced
boundary $D=X\setminus U$, Leray therefore gives
\[
 b_1(X)=2q(X)\le2\sum_x p_g(X_0,x).
\tag{C.DD.8}\label{cdp:eq:C.DD.8}
\]
The resolution is projective, so the equality here
is Hodge decomposition on the actual surface $X$.

Write $c$ for the number of components of $B$,
$b_0=b_0(B)$, and
$\nu=\sum_x(r_x(B)-1)$. Its normalized components
have total genus $g$. Let $r$ be the rank of their
classes in $H_2(E,\Q)$ and put $\varepsilon=c-r$.
The boundary classes on $X$ have exactly
$\varepsilon$ relations. To see the equality, modulo
the independent exceptional classes, the strict
component over $B_i$ is $\pi^*[B_i]/3$.
Indeed its ramification index is three, and the
remaining terms in the divisor pullback are exceptional.
The pullback of $H_2(E,\Q)$ is injective and its
intersection pairing is multiplied by three.
This proves both directions of the asserted relation
count. The span of the boundary classes is
nondegenerate: it contains the positive-square
pullback of $B_{\rm tot}=3\xi$, and the Hodge index
theorem makes the part orthogonal to that class
negative definite.

The pair sequence, exactly as in (C.D4.9)--(C.D4.10),
now gives the version with $b_0$ retained:
\[
\begin{split}
 b_1(U)&=b_1(X)+\varepsilon,\\
 \operatorname{rank}\langle\ ,\ \rangle_U
 &=e(U)+b_0-2+2b_1(X)-b_1(D)+\varepsilon.
\end{split}
\tag{C.DD.9}\label{cdp:eq:C.DD.9}
\]
The second identity follows from
$e(D)=b_0+n_D-b_1(D)$ and the nondegenerate
orthogonal complement of the boundary span. Here
$\langle\ ,\ \rangle_U$ is the compact intersection
pairing, so it is precisely the form preserved by
the normal-cone embedding.

If $\lambda$ is the sum of the first Betti numbers
of the reduced exceptional divisors over the original
surface singularities, normalization of $B$ and the
attaching-branch graph give
\[
 b_1(D)\le2g+\nu-c+b_0+\lambda,
 \qquad e(U)=3(4-2c+2g+\nu).
\tag{C.DD.10}\label{cdp:eq:C.DD.10}
\]
Each exceptional divisor is connected. Its attachments
number $r_x(B)$ at a singular point, giving graph
cycle rank $\nu-c+b_0$. Further boundary blowups
subdivide edges or add trees. The Euler identity
uses the actual unramified degree-three cover
$U\to E\setminus B$ and
$e(B)=2c-2g-\nu$.

\paragraph{No common zero.}
For $v=0$, the divisor $B=R_1+R_2$ is reduced,
has $c=2$, and has arithmetic genus four by
\eqref{cdp:eq:C.DD.2}. Its total delta invariant is therefore
\[
                         \delta(B)=5-g\le4.
\tag{C.DD.11}\label{cdp:eq:C.DD.11}
\]
The local calculation (C.D4.7), or \ref{cdp:st:C.2.RAT.6},
gives $\sum p_g\le1$ and $\lambda\le2$.
The rank $r$ is one or two, so
$\varepsilon=2-r\le1$. Thus $b_1(U)\le3$.
Substitution in \eqref{cdp:eq:C.DD.9}--\eqref{cdp:eq:C.DD.10} gives
\[
                  \operatorname{rank}\langle\ ,\ \rangle_U
                        \ge4g+2\nu-r>0.
\tag{C.DD.12}\label{cdp:eq:C.DD.12}
\]
This calculation includes a disconnected $B$ and
dependent section classes. It contradicts the
normal-cone comparison.

\paragraph{One simple common zero.}
For $v=1$, one has $B=R_1+R_2+F_0$, still reduced.
It is connected because the fibre meets both
horizontal components. Here $c=3$, $b_0=1$,
$\nu\ge2$, and
$\delta(B)=6-g\le5$.
Statement \ref{cdp:st:C.2.RAT.6} supplies
$\sum p_g\le1$ and $\lambda\le2$.
The section and the fibre have independent classes,
so $r=2$ and $\varepsilon=1$.
Equations \eqref{cdp:eq:C.DD.8}--\eqref{cdp:eq:C.DD.10} yield
\[
 b_1(U)\le3,\qquad
 \operatorname{rank}\langle\ ,\ \rangle_U
                         \ge4g+2\nu-6\ge2.
\tag{C.DD.13}\label{cdp:eq:C.DD.13}
\]
This is the same contradiction.

\paragraph{The exact form when the common divisor has degree two.}
Now $N=\OO\oplus\OO(-2)$ by \eqref{cdp:eq:C.DD.4}.
Let $x$ and $y$ be its projections, valued in
$\OO$ and $\OO(-2)$ respectively. The coefficient
of $x^2$ in the irreducible quadratic is a nonzero
constant. Completing its square is a global bundle
automorphism: the mixed coefficient is a section
of $\OO(2)$ and therefore defines an actual map
$\OO(-2)\to\OO$. After a nonzero constant rescaling,
the tensor has the form
\[
 P=a\,y\,(x^2-b y^2),\qquad
 a\in H^0(H,\OO(2)),\quad b\in H^0(H,\OO(4)).
\tag{C.DD.14}\label{cdp:eq:C.DD.14}
\]
Its horizontal root divisors are
$R_1=\{y=0\}$ and $R_2=\{x^2-b y^2=0\}$.
They are disjoint. The quadratic cover has branch
degree at most four. Riemann--Hurwitz and $g\ge1$
force $g=1$, branch degree four, and four distinct
simple zeros of $b$. Thus $R_2$ is already a smooth
elliptic curve. A ruling fibre meets it either at
two transverse points or at a single point of
contact two. This verifies the local models used
in both remaining cases.

If $p^*A=F_1+F_2$ with distinct fibres, the reduced
branch divisor has four components. Its only plane
singularities are ordinary nodes and tangencies of
order two, of types $A_1$ and $A_3$.
Their cubic suspensions have geometric genus zero
and exceptional first Betti number zero, by
(C.RD.10). Consequently $b_1(X)=0$.
The branch divisor is connected, its class rank
is two, and $4\le\nu\le6$.
Thus $\varepsilon=2$, $b_1(D)=\nu-1$, and
\[
 b_1(U)=2,\qquad
 \operatorname{rank}\langle\ ,\ \rangle_U
                       =2\nu-4\ge4.
\tag{C.DD.15}\label{cdp:eq:C.DD.15}
\]
Again the torus comparison is impossible.

\paragraph{A double common zero: normalization of the cyclic cover.}
It remains to take $p^*A=2F_0$. The branch divisor
with multiplicities is $R_1+R_2+2F_0$; its reduced
support $B$ has three components. We use the
normalization of the cyclic cover. Its exact
eigensheaf algebra has underlying module
\[
       \pi_*\OO_{X_0}
       =\OO_E\oplus L^{-1}\oplus L^{-2}(F_0).
\tag{C.DD.16}\label{cdp:eq:C.DD.16}
\]
At the generic point of $F_0$, write the equation
as $w^3=s^2u$ with $u$ a unit. The extra integral
generator is $w^2/s$. This gives the last summand
in \eqref{cdp:eq:C.DD.16}. At the other branch components the
usual generators $1,w,w^2$ apply. These comparisons
extend across the finitely many crossings: the
normal finite surface is Cohen--Macaulay, so its
finite direct image over the smooth surface $E$
is locally free, and each character summand is
the rank-one reflexive extension of its divisorwise
description. This proves \eqref{cdp:eq:C.DD.16} as a sheaf
identity, not only over the nonsingular locus.

Relative duality as in \eqref{cdp:eq:C.DD.7} now gives
\[
 R^1p_*\bigl(L^{-2}(F_0)\bigr)
                       =\det N\otimes\OO(1)=\OO(-1).
\tag{C.DD.17}\label{cdp:eq:C.DD.17}
\]
Thus every summand in \eqref{cdp:eq:C.DD.16} has vanishing
$H^1$, and $H^1(X_0,\OO_{X_0})=0$.
We next prove that all singularities of this
normalized cover are rational.

The only intersections are between the weight-two
branch $F_0$ and a weight-one smooth branch. Their
contact is one or two. The following explicit
resolution works for every contact $k\ge1$.
In coordinates the branch equations are
\[
                F_0=(s=0),\qquad R=(s-u^k=0),
                \qquad w^3=s^2(s-u^k).
\tag{C.DD.18}\label{cdp:eq:C.DD.18}
\]
Blow up their common point until the strict
branches separate, requiring $k$ blowups. At
each step their total branch multiplicity is
$2+1=3$. The exceptional divisor is therefore
unbranched after normalization, and the new
root line and weight-two divisor are exactly
\[
 L'=\sigma^*L(-E_0),\qquad
 F'=\sigma^*F-E_0.
\tag{C.DD.19}\label{cdp:eq:C.DD.19}
\]
For example $s=u s_1$, $w=u w_1$ changes the
equation to
$w_1^3=s_1^2(s_1-u^{k-1})$; the extra normalized
generator changes by
$w^2/s=u(w_1^2/s_1)$.
Both nontrivial eigensheaves consequently satisfy
the exact identities
\[
 (L')^{-1}=\sigma^*(L^{-1})(E_0),\qquad
 (L')^{-2}(F')=\sigma^*(L^{-2}(F))(E_0).
\tag{C.DD.20}\label{cdp:eq:C.DD.20}
\]
For a point blowup of a smooth surface,
$\sigma_*\OO(E_0)=\OO$ and
$R^1\sigma_*\OO(E_0)=0$.
The first follows by extension of functions across
the centre; the second follows from
$0\to\OO\to\OO(E_0)\to\OO_{\PP^1}(-1)\to0$
and the structure-sheaf direct images of the
point blowup. Projection formula applied to
\eqref{cdp:eq:C.DD.20} therefore makes $R^1$ of the finite
cyclic algebra vanish at every step.

For precision, the normalized covers over the
successive blown-up bases have proper birational
maps to $X_0$, induced by their common function
field. Finite direct image and Leray identify the
pushforward of the $R^1$ of each such map with
the just computed $R^1$ of its eigensheaf algebra.
The finite pushforward is faithful, so that $R^1$
is zero. At the end the two branch curves are
smooth and disjoint, and the normalized cover
is smooth: a single branch of multiplicity one
or two is parameterized by $s=t^3$, with $w=t$
or $w=t^2$, up to a unit root. This is a resolution
with $R^1\OO=0$ and proves the claimed rationality.
Its connected exceptional divisors have rational
components and tree incidence graphs: their
$H^1(\OO)$ is a quotient of that $R^1$ and is zero.
Hence their first rational Betti numbers are zero.

In our global double-fibre case it follows that
$b_1(X)=0$ and $\lambda=0$.
The comparison $U\to E\setminus B$ is still an
unramified degree-three cover, since normalization
does not change it off the branch divisor.
Each reduced branch component still has one
strict transform with ramification index three,
so the relation and intersection calculations
\eqref{cdp:eq:C.DD.9}--\eqref{cdp:eq:C.DD.10} apply to this reduced support.
It is connected, $c=3$, its class rank is two,
and $\nu=2$ or $3$. Thus
\[
 b_1(U)=1,\qquad b_1(D)=\nu,\qquad e(U)=3\nu,
 \qquad\operatorname{rank}\langle\ ,\ \rangle_U
                                      =2\nu>0.
\tag{C.DD.21}\label{cdp:eq:C.DD.21}
\]
This is the final contradiction. Every common
divisor allowed by \eqref{cdp:eq:C.DD.4}--\eqref{cdp:eq:C.DD.5} has now
been treated, proving that $C$ cannot have
positive genus.

For the stated irrational ruled consequence,
\ref{cdp:st:A.5.19} leaves a smooth rational or smooth elliptic
image. A rational image would have horizontal
degree two or three. The preceding proof excludes
the first, and \ref{cdp:st:C.2.D4.1} excludes the second.
The elliptic conclusion, its exact minimality
comparison, and its Euler number are those already
proved in \ref{cdp:st:A.5.19}.
\end{proof}

The splitting bounds, all intersection-rank lower
bounds, the exact scalar-zero coordinate example,
and the normalized cyclic-cover blowup substitutions
are retained in \path{code/E_cubic_double.py} and
its saved output. The local delta-five table is
independently computed in \path{code/A_cubic_delta5.py}.


